% Consumo e impuestos verdes · Parte 1: IVA
% Finanzas Públicas, Maestría en Economía Aplicada, Tec de Monterrey
% Esquema aprobado: CLAUDE.md de esta carpeta (2026-10-12). Compilar con: latexmk -pdf consumo_iva.tex
% Versión para repartir: consumo_iva_handout.tex (misma fuente, opción handout).
\documentclass[aspectratio=169,11pt]{beamer}

\usepackage[T1]{fontenc}
\usepackage{lmodern}
\usepackage[spanish,es-nodecimaldot]{babel}
\unaccentedoperators % operadores sin acento: max, min (la opción es-noaccentedoperators no existe en esta versión de babel)
\usepackage{amsmath,amssymb}
\usepackage{booktabs}
\usepackage{tikz}
\usetikzlibrary{babel,arrows.meta,calc,positioning,decorations.pathreplacing,patterns}
\usepackage{pgfplots}
\pgfplotsset{compat=1.18}

% ---------------------------------------------------------------------------
% Tema sobrio (igual al de la sesión del 26 de septiembre)
% ---------------------------------------------------------------------------
\definecolor{cTitulo}{RGB}{0,55,100}
\definecolor{cOferta}{RGB}{31,87,153}
\definecolor{cDemanda}{RGB}{200,100,0}
\definecolor{cImpuesto}{RGB}{175,30,30}
\definecolor{cRecauda}{RGB}{40,130,80}
\definecolor{cGris}{RGB}{110,110,110}

\usetheme{default}
\setbeamertemplate{navigation symbols}{}
\setbeamercolor{frametitle}{fg=cTitulo}
\setbeamercolor{title}{fg=cTitulo}
\setbeamercolor{structure}{fg=cTitulo}
\setbeamercolor{block title}{bg=cTitulo!12,fg=cTitulo}
\setbeamercolor{block body}{bg=cTitulo!4}
\setbeamercolor{block title alerted}{bg=cImpuesto!12,fg=cImpuesto}
\setbeamercolor{block body alerted}{bg=cImpuesto!4}
\setbeamerfont{frametitle}{size=\large,series=\bfseries}
\setbeamertemplate{itemize items}[circle]
\setbeamertemplate{blocks}[rounded]
\setbeamertemplate{footline}{%
  \hfill{\small\color{cGris}Consumo e impuestos verdes · IVA\quad\insertframenumber/\inserttotalframenumber}\hspace*{1em}\vspace*{0.6em}}

% Fuente al pie de cada diapositiva
\newcommand{\fuente}[1]{\par\vfill{\small\color{cGris}#1\par}}
% Resultado final de una derivación, resaltado
\newcommand{\resultado}[1]{\colorbox{cTitulo!12}{$\displaystyle #1$}}

% ---------------------------------------------------------------------------
% Estilo común de gráficas
% ---------------------------------------------------------------------------
\tikzset{
  every node/.append style={font=\small},
  eje/.style={-{Stealth[length=2mm]},thick,black!75},
  flecha/.style={-{Stealth[length=2mm]},thick,cTitulo},
  guia/.style={thin,dashed,black!45},
  punto/.style={circle,fill=black,inner sep=1.2pt},
  caja/.style={draw=cTitulo,rounded corners,fill=cTitulo!5,align=center,inner sep=4pt},
  cajafisco/.style={draw=cRecauda,rounded corners,fill=cRecauda!10,align=center,inner sep=3pt},
}

\title{IVA en México: eficiencia, equidad y lecciones internacionales}
\subtitle{Consumo e impuestos verdes · Parte 1}
\author{Héctor Juan Villarreal}
\institute{Finanzas Públicas · Maestría en Economía Aplicada\\Escuela de Gobierno y Transformación Pública · Tec de Monterrey}
\date{Octubre de 2026}

\begin{document}

% =========================================================================
% I. MOTIVACIÓN
% =========================================================================
\section{Motivación}

% 1 -----------------------------------------------------------------------
\begin{frame}[plain]
  \titlepage
\end{frame}

% 2 -----------------------------------------------------------------------
\begin{frame}{Dos paradigmas}
  \begin{columns}[T]
    \begin{column}{0.48\textwidth}
      \begin{block}{Paradigma viejo}
        Eficiencia contra equidad.\\[3pt]
        En los ochenta y noventa: \emph{recauda con eficiencia y redistribuye con el gasto}. El impuesto no se ocupa de quién es pobre.
      \end{block}
    \end{column}
    \begin{column}{0.48\textwidth}
      \pause
      \begin{block}{Paradigma nuevo}
        La equidad sí importa en el diseño del impuesto.\\[3pt]
        La pregunta deja de ser «¿eficiencia o equidad?» y pasa a ser «¿con qué instrumento logro cada cosa?».
      \end{block}
    \end{column}
  \end{columns}
  \vspace{0.8em}
  {\small Impuestos directos: ISR, predial. Indirectos: IVA, IEPS. Las cuotas de seguridad social son la tercera pieza (sesión del 26 de septiembre).}
  \pause
  \begin{alertblock}{Pregunta}
    ¿Qué cambió?
  \end{alertblock}
\end{frame}

% 3 -----------------------------------------------------------------------
\begin{frame}{México recauda muy poco. ¿Qué pasó?}
  \begin{columns}[T]
    \begin{column}{0.42\textwidth}
      \begin{itemize}
        \item Hay IVA en 175 jurisdicciones. Francia lo aplicó en 1954.
        \item En los países en desarrollo es la principal fuente de recaudación.
        \item México cobra 16\,\% y recauda menos que el promedio de América Latina.
      \end{itemize}
    \end{column}
    \begin{column}{0.56\textwidth}
      \centering
      % Categoría 5110 de la OCDE (impuestos generales sobre bienes y servicios), 2021,
      % base Revenue Statistics in LAC consultada el 12-10-2026. El gráfico 1 de Barreix et al.
      % traía México 4.3, Uruguay 7.9, Brasil 12.5, ALC 6.4: no coinciden con la base actual.
      \begin{tikzpicture}
        \begin{axis}[
          width=8.2cm, height=4.3cm,
          ybar, bar width=14pt, bar shift=0pt,
          ymin=0, ymax=14.5,
          ylabel={Impuestos generales al consumo,\\\% del PIB (2021)},
          ylabel style={align=center},
          ylabel style={font=\small},
          symbolic x coords={Brasil,Chile,Uruguay,OCDE,ALC,Colombia,México},
          xtick={Brasil,Chile,Uruguay,OCDE,ALC,Colombia,México},
          x tick label style={font=\small,rotate=35,anchor=north east,inner sep=1pt},
          y tick label style={font=\small},
          nodes near coords,
          every node near coord/.append style={font=\small},
          axis lines*=left,
          enlarge x limits=0.09,
        ]
          \addplot[fill=cOferta!55,draw=cOferta] coordinates {(Brasil,12.1) (Chile,9.5) (Uruguay,7.0) (Colombia,6.7)};
          \addplot[fill=cGris!40,draw=cGris] coordinates {(OCDE,7.1) (ALC,6.6)};
          \addplot[fill=cImpuesto!60,draw=cImpuesto] coordinates {(México,4.2)};
        \end{axis}
      \end{tikzpicture}
    \end{column}
  \end{columns}
  \pause
  \begin{alertblock}{Pregunta}
    Con una tasa de 16\,\%, ¿por qué recaudamos tan poco?
  \end{alertblock}
  \fuente{Categoría 5110; Brasil suma ICMS, PIS/COFINS e ISS. Base OCDE-BID-CEPAL-CIAT de \emph{Estadísticas tributarias en ALC} (octubre de 2026). 175 jurisdicciones: OCDE (2024), citada por Barreix et al.}
\end{frame}

% 4 -----------------------------------------------------------------------
\begin{frame}{Objetivo de hoy}
  \begin{block}{Tender un puente}
    De la evidencia internacional sobre eficiencia (FMI) a las propuestas de equidad de América Latina (Barreix), y de ahí a la reforma fiscal en México.
  \end{block}
  \vspace{0.6em}
  \centering
  \begin{tikzpicture}[node distance=4mm]
    \node[caja,text width=2.2cm] (a) {Mecánica del IVA};
    \node[caja,text width=2.2cm,right=of a] (b) {No gravar insumos};
    \node[caja,text width=2.2cm,right=of b] (c) {¿Tasa uniforme?};
    \node[caja,text width=2.2cm,right=of c] (d) {Incidencia y ganancias};
    \node[caja,text width=2.2cm,right=of d] (e) {Barreix y México};
    \draw[flecha] (a) -- (b);
    \draw[flecha] (b) -- (c);
    \draw[flecha] (c) -- (d);
    \draw[flecha] (d) -- (e);
    \node[below=3mm of b,text width=2.4cm,align=center,color=cGris] {FMI};
    \node[below=3mm of c,text width=2.4cm,align=center,color=cGris] {FMI y Salanié\\(entra al parcial)};
    \node[below=3mm of e,text width=2.4cm,align=center,color=cGris] {Barreix et al.};
  \end{tikzpicture}
  \fuente{de Mooij, Hebous y Keen (2025), IMF WP/25/165; Barreix, Rocha y Tostes Neto (2025), CIATBlog.}
\end{frame}

% =========================================================================
% II. MECÁNICA
% =========================================================================
\section{Mecánica}

% 5 -----------------------------------------------------------------------
\begin{frame}{IVA e impuesto a las ventas: los dos gravan el consumo, solo que…}
  \centering
  \begin{tikzpicture}[x=1cm,y=1cm]
    \node[caja,text width=2.3cm] (p) at (0,0) {Proveedor de insumos\\vende 100};
    \node[caja,text width=2.3cm] (f) at (3.8,0) {Fabricante\\vende 300};
    \node[caja,text width=2.3cm] (m) at (7.6,0) {Minorista\\vende 500};
    \node[caja,text width=2.3cm] (c) at (11.4,0) {Consumidor\\paga 580};
    \draw[flecha] (p) -- (f); \draw[flecha] (f) -- (m); \draw[flecha] (m) -- (c);
    \onslide<2->{
    \node[cajafisco,text width=2.6cm] at (0,-1.6) {IVA cobrado 16\\acredita 0\\\textbf{entera 16}};
    \node[cajafisco,text width=2.6cm] at (3.8,-1.6) {IVA cobrado 48\\acredita 16\\\textbf{entera 32}};
    \node[cajafisco,text width=2.6cm] at (7.6,-1.6) {IVA cobrado 80\\acredita 48\\\textbf{entera 32}};
    \node[align=center,text width=2.8cm,color=cRecauda] at (11.4,-1.6) {Total: $16+32+32$\\$= 80 = 0.16\times 500$};}
  \end{tikzpicture}
  \par\medskip
  \begin{itemize}
    \onslide<3->{\item Cada empresa paga débito menos crédito. Lo recaudado es $t$ por la venta final: igual que un impuesto a las ventas al menudeo.}
    \onslide<4->{\item La diferencia es el \textbf{cobro fraccionado}: si el minorista evade, el fisco ya cobró 48 de los 80. Y cada factura deja huella documental.}
    \onslide<5->{\item Un impuesto a los ingresos brutos cobra en cada venta sin acreditar: el impuesto se acumula en cascada.}
  \end{itemize}
  \fuente{Tasa de 16\,\%. de Mooij, Hebous y Keen (2025), sec.~2.3.}
\end{frame}

% 6 -----------------------------------------------------------------------
\begin{frame}{Principio de destino: se grava donde se consume}
  \begin{columns}[T]
    \begin{column}{0.52\textwidth}
      \begin{itemize}
        \item \textbf{Exportaciones a tasa cero}: salen sin IVA y la empresa recupera el IVA de sus insumos.
        \item \textbf{Importaciones gravadas}: pagan la misma tasa que el producto nacional.
        \item Resultado: el IVA no distingue entre lo hecho aquí y lo hecho afuera. Es neutral al comercio.
      \end{itemize}
      \pause
      \begin{block}{«El impuesto interno del libre comercio»}
        Por eso fue condición para entrar a la Unión Europea. Un arancel, en cambio, grava la producción y viola la eficiencia productiva.
      \end{block}
    \end{column}
    \begin{column}{0.44\textwidth}
      \pause
      \begin{block}{Brasil, hoy}
        Su ICMS estatal se cobra en buena parte en origen: el estado productor se queda con la recaudación.\\[3pt]
        La reforma lleva el nuevo IBS al principio de destino, con una transición de 50 años.
      \end{block}
    \end{column}
  \end{columns}
  \fuente{de Mooij, Hebous y Keen (2025), secs.~2.1 y 2.3; Barreix, Rocha y Tostes Neto (2025).}
\end{frame}

% 7 -----------------------------------------------------------------------
\begin{frame}{Tasa cero no es lo mismo que exención}
  \begin{columns}[T]
    \begin{column}{0.44\textwidth}
      \centering
      \begin{tabular}{lcc}
        \toprule
        & Cobra IVA & Acredita\\
        \midrule
        Tasa general & sí & sí\\
        Tasa cero & no & sí\\
        Exención & no & \textcolor{cImpuesto}{\textbf{no}}\\
        \bottomrule
      \end{tabular}
      \par\medskip
      \raggedright
      Con la exención, el IVA de los insumos se queda pegado: se vuelve un impuesto a los insumos.
      \par\medskip
      {\small México: alimentos y medicinas a tasa cero (LIVA, art.~2-A). Exenciones: arts.~9 (enajenación), 15 (servicios) y 20 (uso o goce temporal).}
    \end{column}
    \begin{column}{0.54\textwidth}
      \pause
      \begin{block}{Exentamos al fabricante de la diapositiva 5}
        \small
        \begin{itemize}
          \item No recupera los 16 de IVA de sus insumos: vende en 316, sin IVA.
          \item El minorista no tiene IVA que acreditar: vende en 516 y cobra $0.16\times516 = 82.56$.
          \item Recaudación: $16 + 0 + 82.56 = 98.56$, no 80.
          \item El consumidor paga 598.56, no 580.
        \end{itemize}
      \end{block}
      \pause
      Exentar una etapa intermedia \textbf{sube} la recaudación y el precio final: el impuesto se cobra dos veces sobre lo mismo.
    \end{column}
  \end{columns}
  \fuente{Ley del Impuesto al Valor Agregado, arts.~2-A, 9, 15 y 20. Ejemplo propio; supone que el fabricante traslada el IVA no acreditable.}
\end{frame}

% =========================================================================
% III. EFICIENCIA EN LA PRODUCCIÓN
% =========================================================================
\section{Eficiencia en la producción}

% 8 -----------------------------------------------------------------------
\begin{frame}{No gravar insumos: Diamond y Mirrlees (1971)}
  \begin{itemize}
    \item Gravar el consumo y gravar la producción distorsionan los dos. Uno pensaría que conviene usar un poco de cada uno.
    \pause
    \item No es así: si el gobierno puede gravar el consumo final a las tasas que quiera, conviene gravar \textbf{solo} el consumo.
    \pause
    \item Gravar insumos deja la economía dentro de su frontera de posibilidades: se produce menos de lo que se podría. Y esa producción extra tendría algún buen uso.
  \end{itemize}
  \pause
  \begin{block}{Eficiencia productiva}
    El sistema fiscal no debe alterar cómo se organiza la producción: qué insumos se usan ni quién produce qué.
  \end{block}
  \pause
  {\small Supuestos: competencia; rentas (utilidades extraordinarias) gravables a la tasa que se quiera; ninguna restricción para gravar el consumo.}
  \fuente{Diamond y Mirrlees (1971), \emph{AER} 61(1): 8--27 y 61(3): 261--278; de Mooij, Hebous y Keen (2025), sec.~2.1 y nota 4; Salanié, sec.~3.2.}
\end{frame}

% 9 -----------------------------------------------------------------------
\begin{frame}{Cuánto cuesta gravar insumos}
  \begin{columns}[T]
    \begin{column}{0.47\textwidth}
      \[
        \frac{CE}{\text{costo}} = \tfrac12\,\alpha\,(1-\alpha)\,\sigma\,t^2
      \]
      {\small $\alpha$: participación de insumos gravados; $\sigma$: sustitución entre insumos; $t$: tasa ad valorem.}
      \pause
      \begin{itemize}
        \item Pariente de $\tfrac12\,t^2\,\varepsilon$ (26 de septiembre, diap.~27).
        \item $\alpha(1-\alpha)$: pérdida chica si se grava casi ningún insumo o casi todos (los precios relativos no cambian).
      \end{itemize}
      \pause
      \begin{block}{Ejemplo del FMI}
        Con $\sigma=2$, $\alpha=0.5$ y $t=5$\,\%: 0.06\,\% del costo y 0.82\,\% de lo recaudado.
      \end{block}
    \end{column}
    \begin{column}{0.5\textwidth}
      \centering
      \begin{tikzpicture}
        \begin{axis}[
          width=7.2cm, height=5.2cm,
          xmin=0, xmax=1, ymin=0, ymax=0.075,
          xlabel={$\alpha$ (con $t=5$\,\%)},
          ylabel={$CE$, \% del costo},
          xlabel style={font=\small}, ylabel style={font=\small},
          tick label style={font=\small},
          yticklabel style={/pgf/number format/fixed,/pgf/number format/precision=2},
          scaled y ticks=false,
          axis lines*=left,
          domain=0:1, samples=60,
        ]
          \addplot[very thick,cImpuesto] {100*0.5*x*(1-x)*2*0.05^2};
          \addplot[very thick,cOferta,dash pattern=on 5pt off 2pt] {100*0.5*x*(1-x)*1*0.05^2};
          \node[cImpuesto,font=\small] at (axis cs:0.5,0.069) {$\sigma=2$};
          \node[cOferta,font=\small] at (axis cs:0.5,0.025) {$\sigma=1$};
          \node[punto] at (axis cs:0.5,0.0625) {};
        \end{axis}
      \end{tikzpicture}
    \end{column}
  \end{columns}
  \fuente{Keen (2014), \emph{Indian Growth and Development Review} 7(2): 181--201; FMI, fig.~1 (redibujada).}
\end{frame}

% 10 ----------------------------------------------------------------------
\begin{frame}{La cascada no es, por sí sola, ineficiente}
  \begin{itemize}
    \item Un impuesto de 5\,\% sobre un insumo, y otra vez sobre lo que se produce con él: $1.05\times1.05 = 1.1025$. La tasa efectiva es 10.25\,\%.
    \pause
    \item Si cada unidad de producto requiere una cantidad fija de insumo (coeficientes técnicos fijos), nadie cambia cómo produce: \textbf{no hay carga excesiva}.
    \pause
    \item El costo viene de otro lado:
      \begin{itemize}
        \item de la \textbf{sustitución}: la empresa usa menos del insumo gravado aunque sea el mejor;
        \item de la \textbf{integración vertical} para brincarse etapas, que cuesta;
        \item de la \textbf{opacidad}: nadie sabe cuál es la tasa efectiva sobre cada bien.
      \end{itemize}
  \end{itemize}
  \pause
  \begin{block}{Contraste con Barreix et al.}
    Barreix et al. llaman «efecto cascada» a lo que dejan las exenciones y lo tratan como ineficiencia. El FMI matiza: la cascada importa en la medida en que hay sustitución, integración u opacidad.
  \end{block}
  \fuente{de Mooij, Hebous y Keen (2025), sec.~2.1, notas 7 y 8; Barreix, Rocha y Tostes Neto (2025).}
\end{frame}

% 11 ----------------------------------------------------------------------
\begin{frame}{Las exenciones deshacen el IVA}
  \begin{itemize}
    \item Una exención es un impuesto a los insumos: trae todas las distorsiones de la cascada.
    \item Empuja al \textbf{autosuministro}: la empresa hace en casa lo que compraría, para no cargar IVA que no puede acreditar.
    \item Vuelve opacas las tasas efectivas.
    \item Las importaciones llegan con tasa cero en origen y sin IVA en sus insumos: ganan ventaja frente al producto nacional exento.
  \end{itemize}
  \pause
  \begin{block}{Política industrial vía IVA}
    \begin{itemize}
      \item Tasa reducida a un insumo que compran contribuyentes registrados: no tiene efecto, porque el comprador acredita todo.
      \item Exención a ese insumo: contraproducente.
    \end{itemize}
  \end{block}
  \pause
  \begin{alertblock}{Pregunta}
    Si exentamos un insumo que compran empresas registradas, ¿el precio final baja o sube?
  \end{alertblock}
  \fuente{de Mooij, Hebous y Keen (2025), secs.~2.3 y 3.3.}
\end{frame}

% =========================================================================
% IV. ¿TASA UNIFORME?
% =========================================================================
\section{¿Tasa uniforme?}

% 12 ----------------------------------------------------------------------
\begin{frame}{Un IVA distorsiona en dos márgenes}
  \begin{columns}[T]
    \begin{column}{0.5\textwidth}
      \begin{block}{1. Entre bienes}
        Si los bienes pagan tasas distintas, los precios relativos cambian: la tortilla con tasa cero frente al refresco con 16\,\%.
      \end{block}
      \pause
      \begin{block}{2. Entre consumo y ocio}
        El ocio no se puede gravar. Con IVA, conviene más pintar la casa uno mismo que trabajar horas extra para pagarle al pintor.
      \end{block}
    \end{column}
    \begin{column}{0.46\textwidth}
      \pause
      \begin{alertblock}{Un IVA uniforme es un impuesto al trabajo}
        Grava todo lo que el salario compra, igual que un impuesto al salario. Lo vimos con Harberger: sesión del 26 de septiembre, diap.~19.
      \end{alertblock}
      \pause
      Por eso la tasa uniforme no elimina toda distorsión. Quita la primera y deja la segunda.
    \end{column}
  \end{columns}
  \fuente{de Mooij, Hebous y Keen (2025), sec.~2.2.}
\end{frame}

% 13 ----------------------------------------------------------------------
\begin{frame}{La regla de Ramsey}
  Precios al productor fijos y normalizados a 1: $q_i = 1 + t_i$. El gobierno necesita recaudar $R$.
  \begin{align*}
    \onslide<2->{\text{Problema:}\quad & \max_{t}\ V(q, I) \quad\text{s.a.}\quad \textstyle\sum_k t_k\, x_k(q,I) \ge R\\}
    \onslide<3->{\text{Condición para } t_i\text{ (Roy):}\quad & -\alpha\, x_i + \lambda\Big(x_i + \textstyle\sum_k t_k\, \frac{\partial x_k}{\partial q_i}\Big) = 0\\}
    \onslide<4->{\text{Slutsky y simetría:}\quad & \tfrac{\partial x_k}{\partial q_i} = S_{ik} - x_i\,\tfrac{\partial x_k}{\partial I}, \qquad S_{ik} = S_{ki}\\}
    \onslide<5->{\text{Resultado:}\quad & \resultado{\frac{\sum_k t_k\, S_{ik}}{x_i} = -\theta \quad \text{para todo } i}}
  \end{align*}
  \onslide<6->{
  \begin{block}{Lectura}
    $\sum_k t_k S_{ik}$ es, en primera aproximación, cuánto cae la demanda \emph{compensada} de $i$ por todos los impuestos. La regla pide que caiga \textbf{en la misma proporción} para todos los bienes. No pide tasas iguales.
  \end{block}}
  \fuente{Salanié, cap.~3, sec.~3.1. $\alpha$: utilidad marginal del ingreso; $\theta>0$ no depende de $i$.}
\end{frame}

% 14 ----------------------------------------------------------------------
\begin{frame}{La elasticidad inversa es un caso particular}
  \begin{columns}[T]
    \begin{column}{0.5\textwidth}
      Si las demandas compensadas son independientes ($S_{ik}=0$ para $i\neq k$):
      \[
        t_i\,\frac{S_{ii}}{x_i} = -\theta
        \quad\Longrightarrow\quad
        \resultado{\frac{t_i}{1+t_i} = \frac{\theta}{e_i}}
      \]
      con $e_i$ la elasticidad compensada de $i$, en valor absoluto.
      \pause
      \begin{itemize}
        \item Más tasa a lo que tiene demanda inelástica.
        \item Lectura del FMI: demanda inelástica quiere decir que el bien \textbf{acompaña al ocio}. Si sube su precio y no lo dejo, trabajo más para pagarlo.
      \end{itemize}
    \end{column}
    \begin{column}{0.46\textwidth}
      \pause
      \begin{block}{Corlett y Hague (1953)}
        Gravar menos lo que acompaña al trabajo: guarderías, transporte al trabajo.\\[3pt]
        Gravar más lo que acompaña al ocio: películas en línea, quizá alcohol.
      \end{block}
      \pause
      \begin{alertblock}{Pregunta}
        ¿Qué bien gravarían más y cuál menos? ¿Por qué?
      \end{alertblock}
    \end{column}
  \end{columns}
  \fuente{Salanié, sec.~3.1.3; de Mooij, Hebous y Keen (2025), sec.~2.2; Corlett y Hague (1953), \emph{RES} 21(1): 21--30.}
\end{frame}

% 15 ----------------------------------------------------------------------
\begin{frame}{Atkinson y Stiglitz: con buen ISR, basta una tasa}
  \begin{block}{Resultado}
    Si hay un impuesto óptimo al ingreso y las preferencias son \textbf{débilmente separables} entre consumo y ocio, la tasa uniforme es óptima: los impuestos indirectos no agregan nada.
  \end{block}
  \pause
  \begin{itemize}
    \item Es la versión formal del paradigma viejo de la diapositiva 2: recauda con el IVA, redistribuye con el ISR y el gasto.
    \item Lo que justifica una tasa distinta es la relación de cada bien con el ocio, no quién lo consume.
  \end{itemize}
  \pause
  \begin{block}{Lo que dice el FMI}
    \begin{itemize}
      \item No hay teorema que haga óptima la uniformidad siempre, pero apartarse de ella exige razones contundentes.
      \item Cuanto más ricos los instrumentos para redistribuir, más débil el caso para diferenciar tasas.
    \end{itemize}
  \end{block}
  \fuente{Atkinson y Stiglitz (1976), \emph{JPubE} 6(1--2): 55--75; Salanié, sec.~5.2; de Mooij, Hebous y Keen (2025), sec.~2.2.}
\end{frame}

% 16 ----------------------------------------------------------------------
\begin{frame}{Redistribuir con tasas reducidas rinde poco}
  \begin{columns}[T]
    \begin{column}{0.48\textwidth}
      \begin{block}{Los ricos también comen}
        La tasa cero a la tortilla ayuda al hogar pobre, pero el hogar rico compra más alimentos en pesos y se lleva más subsidio.
      \end{block}
      \pause
      Lo que cuenta es qué tan \textbf{distintas} son las canastas de ricos y pobres. Si se parecen, diferenciar tasas redistribuye muy poco (Sah, 1983).
    \end{column}
    \begin{column}{0.48\textwidth}
      \pause
      \begin{block}{México, 2026}
        \begin{itemize}
          \item Tasa cero: 1.5\,\% del PIB (alimentos, 1.26\,\%).
          \item Exenciones: 0.21\,\% del PIB.
          \item La mitad más rica (deciles VI a X) se lleva 58.5\,\% de la tasa cero en alimentos.
        \end{itemize}
      \end{block}
      \pause
      Como proporción de su ingreso, pesa más en los deciles bajos. Es progresiva, pero cara.
    \end{column}
  \end{columns}
  \fuente{SHCP (2026), \emph{Renuncias Recaudatorias 2026}, cuadros 3.3.1--3.3.2, gráficas 28--29; Sah (1983), \emph{JPubE} 20(1): 89--101.}
\end{frame}

% =========================================================================
% V. INCIDENCIA Y GANANCIAS
% =========================================================================
\section{Incidencia y ganancias}

% 17 ----------------------------------------------------------------------
\begin{frame}{¿Quién paga el IVA?}
  \begin{columns}[T]
    \begin{column}{0.54\textwidth}
      \begin{itemize}
        \item \textbf{Tasa general}: traslado casi completo (zona euro, 1999--2013; Alemania).
        \item \textbf{Tasas reducidas}: alrededor de 30\,\% en la zona euro.
        \item \textbf{Restaurantes} en Francia y Suecia: 14--25\,\%.
      \end{itemize}
      \pause
      \begin{block}{Lectura con la sesión anterior}
        Traslado $= \frac{e_S}{e_S+e_D}$ (diap.~8). Base amplia: poco a dónde sustituir, $e_D$ chica, traslado alto. Tasa a un solo sector: el consumidor sí sustituye.
      \end{block}
    \end{column}
    \begin{column}{0.42\textwidth}
      \pause
      \begin{alertblock}{¿Y en países en desarrollo?}
        Casi no hay evidencia.\\[3pt]
        La informalidad podría reducir el traslado: es fácil salirse del IVA (sesión del 26 de septiembre, diap.~14).\\[3pt]
        Con una tasa cero que no se traslada, el subsidio se queda en el vendedor.
      \end{alertblock}
    \end{column}
  \end{columns}
  \fuente{de Mooij, Hebous y Keen (2025), sec.~2.5; Benedek et al. (2020), \emph{ITPF} 27(4): 890--930.}
\end{frame}

% 18 ----------------------------------------------------------------------
\begin{frame}{Cuánto se gana con un IVA}
  \[
    \text{Ganancia} \approx \Delta\,\text{CMFP} \times \text{recaudación inicial}
  \]
  {\small CMFP: costo marginal de los fondos públicos.}
  \begin{align*}
    \onslide<2->{\text{Auriol y Warlters (2012):}\ & \text{mediana } 1.23 \text{ sin IVA},\ 1.17 \text{ con IVA}\ \Rightarrow\ \Delta = 0.06\\}
    \onslide<3->{\text{Recaudación de 12 a 15 \% del PIB:}\ & 0.06\times 12 = 0.72 \ \text{a}\ 0.06\times15 = 0.9 \ \text{(\% del PIB)}}
  \end{align*}
  \onslide<4->{
  \begin{itemize}
    \item Es una comparación entre 38 países africanos, no una caída medida en el tiempo. El FMI lo llama un supuesto «heroico».
    \item Con un segundo método del FMI, la ganancia va de 0.3 a 2.1\,\% del PIB.
    \item No es una ganancia de una vez: se repite cada año.
  \end{itemize}}
  \onslide<5->{
  \begin{block}{Enlace con la sesión anterior}
    Cada peso adicional cuesta $1+\frac{\eta}{1-\eta}$ (diap.~34). El IVA ayuda porque su base responde poco a la tasa: baja $\eta$.
  \end{block}}
  \fuente{de Mooij, Hebous y Keen (2025), sec.~2.4; Auriol y Warlters (2012), \emph{JDE} 97(1): 58--72.}
\end{frame}

% =========================================================================
% VI. BARREIX
% =========================================================================
\section{Barreix: el puente hacia la equidad}

% 19 ----------------------------------------------------------------------
\begin{frame}{IVA personalizado: cobrar parejo y devolver a quien lo necesita}
  \begin{columns}[T]
    \begin{column}{0.5\textwidth}
      \begin{itemize}
        \item Base amplia y tasa uniforme.
        \item El fisco devuelve a los hogares pobres el IVA que pagan en una canasta típica.
        \item Se paga por transferencia electrónica, con los padrones de los programas sociales.
      \end{itemize}
      \pause
      \begin{block}{Es Atkinson y Stiglitz en la práctica}
        El IVA recauda; la devolución focalizada redistribuye. El FMI lo cita como alternativa a las tasas reducidas (nota 28).
      \end{block}
    \end{column}
    \begin{column}{0.46\textwidth}
      \pause
      \begin{block}{En América Latina}
        \small
        \begin{itemize}
          \item Uruguay (2012): hogares vulnerables, si pagan con tarjeta.
          \item Ecuador (2017): adultos mayores pobres y personas con discapacidad.
          \item Colombia (2020): compensación a hogares pobres.
          \item Bolivia (2021): según las facturas a nombre de la persona.
          \item Argentina (2016): 15\,\% con tarjeta de débito a jubilados y a la AUH; efecto muy limitado.
        \end{itemize}
      \end{block}
    \end{column}
  \end{columns}
  \fuente{Barreix et al. (2022), «El IVA personalizado revisado», BID, IDB-DP-939 (doi 10.18235/0004147), pp.~20--21; de Mooij, Hebous y Keen (2025), nota 28.}
\end{frame}

% 20 ----------------------------------------------------------------------
\begin{frame}{Brasil: tres innovaciones}
  \small
  \begin{columns}[T]
    \begin{column}{0.36\textwidth}
      \begin{block}{1. IVA dual}
        CBS federal; IBS estatal y municipal. Un comité gestor recauda el IBS y lo reparte por destino. Transición de 50 años.
      \end{block}
      \begin{block}{3. Cashback}
        Hogares con hasta medio salario mínimo per cápita: 100\,\% de la CBS y 20\,\% del IBS en gas, luz, agua, drenaje y telecom.; 20\,\% en lo demás. Antes: Devolve ICMS (Rio Grande do Sul, 2021).
      \end{block}
    \end{column}
    \begin{column}{0.62\textwidth}
      \begin{block}{2. Pago fraccionado (\emph{split payment})}
        \centering
        \begin{tikzpicture}[x=1cm,y=1cm,font=\small]
          \node[caja,text width=1.9cm] (prov) at (0,0) {Proveedor\\vende 100};
          \node[caja,text width=1.9cm] (comp) at (6.2,0) {Comprador\\paga 126.5};
          \node[caja,text width=2.1cm] (ent) at (3.1,-1.5) {Entidad\\de pago};
          \node[cajafisco,text width=2.4cm] (fis) at (3.1,-3.3) {Fisco y comité gestor};
          \draw[flecha] (prov) -- node[above] {100 + 26.5 de IBS/CBS} (comp);
          \draw[flecha] (comp) |- (ent);
          \draw[flecha,cRecauda] ([xshift=-4mm]ent.south) -- node[left,align=right,font=\small] {¿tiene\\créditos?} ([xshift=-4mm]fis.north);
          \draw[flecha,cRecauda] ([xshift=4mm]ent.south) -- node[right,align=left,font=\small] {si no:\\26.5} ([xshift=4mm]fis.north);
          \draw[flecha] (ent.west) -| node[pos=0.3,above] {100} (prov.south);
        \end{tikzpicture}
        \par
        Si el proveedor no tiene créditos, la entidad de pago separa los 26.5 y le pasa 100. El 26.5\,\% es una tasa estimada, no fijada en ley.
      \end{block}
    \end{column}
  \end{columns}
  \fuente{Barreix, Rocha y Tostes Neto (2025), gráfico 3 (redibujado); LC 214/2025, arts.~31--35, 113 y 118.}
\end{frame}

% 21 ----------------------------------------------------------------------
\begin{frame}{FMI y Barreix: ¿visiones complementarias?}
  \begin{columns}[T]
    \begin{column}{0.46\textwidth}
      \begin{block}{En qué coinciden}
        \begin{itemize}
          \item Tasa uniforme.
          \item Nada de exenciones amplias.
        \end{itemize}
        El FMI pone la eficiencia; Barreix pone la equidad, vía devoluciones.
      \end{block}
    \end{column}
    \begin{column}{0.5\textwidth}
      \pause
      \begin{alertblock}{Lectura crítica}
        \begin{enumerate}
          \item La reforma brasileña no es uniforme: tasas reducidas en 30\,\% y 60\,\% y tasa cero, Zona Franca de Manaos, Simples Nacional. ¿El cashback y las tasas reducidas se duplican?
          \pause
          \item Barreix et al. (cuadro 1): la informalidad vuelve «invisibles» a los deciles bajos. Pero el cashback también necesita identificarlos.
        \end{enumerate}
      \end{alertblock}
    \end{column}
  \end{columns}
  \fuente{Barreix, Rocha y Tostes Neto (2025) y nota larga de los mismos autores (mayo de 2025), cuadro 1; de Mooij, Hebous y Keen (2025), sec.~2.2.}
\end{frame}

% =========================================================================
% VII. MÉXICO
% =========================================================================
\section{México}

% 22 ----------------------------------------------------------------------
\begin{frame}{La oportunidad para México}
  \begin{itemize}
    \item \textbf{Base erosionada}: tasa cero en alimentos y medicinas, exenciones, tasa de 8\,\% en la región fronteriza (estímulo vigente en 2026).
    \item \textbf{Informalidad}: buena parte del consumo pasa por vendedores que no cobran IVA.
    \item \textbf{Recaudación baja frente a pares} (diapositiva 3).
  \end{itemize}
  \pause
  \begin{block}{Cuánto deja de entrar}
    Tasa cero y exenciones del IVA: 638.7 mil millones de pesos en 2026, 1.7\,\% del PIB.\\[2pt]
    {\small No todo se recaudaría: con el IVA a los alimentos, parte del consumo se iría a la informalidad.}
  \end{block}
  \pause
  \begin{block}{Lo que ya existe}
    Factura electrónica (CFDI) en toda la cadena y transferencias sociales pagadas a cuentas bancarias.
  \end{block}
  \fuente{SHCP (2026), \emph{Renuncias Recaudatorias 2026}, cuadros 3.3.1 y 3.3.2 (suma propia); decreto de estímulos fiscales en las regiones fronterizas norte y sur (DOF, 31-12-2025).}
\end{frame}

% 23 ----------------------------------------------------------------------
\begin{frame}{Riesgos y desafíos}
  \begin{itemize}
    \item \textbf{Padrones}: devolver exige saber a quién. ¿Los padrones de los programas sociales sirven para eso?
    \item \textbf{Política}: gravar alimentos y medicinas es muy impopular, aunque se devuelva más de lo que se cobra.
    \item \textbf{Liquidez}: el hogar paga el IVA hoy y recibe la devolución después. Si tarda, el pobre financia al fisco.
    \item \textbf{Confianza}: si la gente no ve llegar la devolución, la reforma se cae.
  \end{itemize}
  \pause
  \begin{alertblock}{La informalidad del lado del hogar}
    El hogar que vive del trabajo informal no aparece en el SAT ni siempre en los padrones. ¿A quién se le devuelve?
  \end{alertblock}
\end{frame}

% 24 ----------------------------------------------------------------------
\begin{frame}{Para la reforma fiscal: ¿qué harían?}
  Propuestas para discutir, no conclusiones:
  \begin{enumerate}
    \item Avanzar a la tasa única con un plan gradual.
    \item Un piloto de IVA personalizado en un grupo bien identificado.
    \item Más administración tributaria y más factura electrónica.
    \item Explicar con claridad cómo se protege a los hogares vulnerables, antes de cobrar.
  \end{enumerate}
  \pause
  \begin{alertblock}{Pregunta}
    ¿Por dónde empezarían y qué dejarían para después?
  \end{alertblock}
\end{frame}

% =========================================================================
% VIII. PUENTE
% =========================================================================
\section{Puente}

% 25 ----------------------------------------------------------------------
\begin{frame}{El IVA es eficiente, pero en México se percibe injusto}
  \begin{itemize}
    \item Un IVA amplio y uniforme, con devoluciones focalizadas, puede ser eficiente y justo a la vez.
    \pause
    \item El FMI pone el marco de eficiencia; Barreix, el puente hacia la equidad.
  \end{itemize}
  \pause
  \begin{alertblock}{Recaudar no es corregir}
    Si el objetivo es corregir un daño ---el CO$_2$ de la gasolina, el azúcar del refresco---, la herramienta no es una tasa distinta de IVA. Es un impuesto específico.
  \end{alertblock}
  \pause
  \centering\large Después del descanso: impuestos verdes.
  \fuente{de Mooij, Hebous y Keen (2025), sec.~3.3.}
\end{frame}

\end{document}
